\documentclass[]{article}

\usepackage{amsmath}
\usepackage{multirow}
\usepackage{natbib}
\usepackage{graphicx} 


\begin{document}

\subsection{Thermodynamic stability classification}

The first stage of our pipeline addresses the prediction of thermodynamic stability, a binary classification task where stable compounds are defined as having an energy above the convex hull of 0 eV/atom. Given the significant class imbalance in our dataset, with only 13.1\% of the 1283 compounds being stable, we employed a Gradient Boosting Classifier (GBC) evaluated with a rigorous Leave-One-Cluster-Out (LOCO) cross-validation. By partitioning the data into 53 clusters based on the A-site element, this strategy ensures the model is tested on its ability to generalize to entirely new chemical families, preventing overly optimistic performance estimates that can arise from random data splits.

The model achieved an Area Under the Precision-Recall Curve (AUC-PR) of 0.353 and an F1-score of 0.320. As illustrated in Figure \ref{fig:image0}, the precision-recall curve demonstrates that the model's performance is substantially better than the no-skill baseline of 0.131 (the positive class frequency). This indicates a strong capability to enrich a pool of candidates with thermodynamically stable materials, which is the primary goal of this initial filtering stage. The chemical heterogeneity of the dataset, where stability rates vary drastically between A-site groups, underscores the necessity of the LOCO approach for building a truly generalizable model that learns the underlying physical principles of stability rather than memorizing compositional patterns.

\begin{figure}[h!]
    \centering
    \includegraphics[width=0.5\textwidth]{../input_files/plots/step_2_pr_curve_stability_1_1776245452.png}
    \caption{Precision-Recall curve for the Gradient Boosting Classifier (GBC) predicting the thermodynamic stability of ABO$_3$ perovskites. The model was evaluated using a rigorous Leave-One-Cluster-Out (LOCO) cross-validation strategy, grouped by the A-site element. The resulting Area Under the Curve (AUC-PR) of 0.353 demonstrates a substantial performance improvement over the no-skill baseline of 0.131, which corresponds to the 13.1\% frequency of stable compounds in the dataset. This illustrates the model's ability to effectively enrich a candidate pool with thermodynamically stable materials far beyond random selection.}
    \label{fig:image0}
\end{figure}

\subsection{Mechanical viability classification and uncertainty quantification}

A central challenge in high-throughput materials screening is the prevalence of sparse and unreliable data for mechanical properties. Our dataset reflects this issue, with elastic constants available for only 16.8\% of compounds and numerous entries containing unphysical values (e.g., negative or extremely large moduli) indicative of calculation failures. To circumvent the pitfalls of regressing on this flawed data, we reframed the problem as a binary classification task.

A Mechanical Viability Classifier (a GBC model) was trained to distinguish physically consistent materials (207 samples with $0 < K_{VRH} < 300$ GPa and $G_{VRH} > 0$) from the 8 unphysical or unstable ones. Evaluated with 5-fold stratified cross-validation, this classifier demonstrated exceptional performance in identifying sound candidates, achieving an accuracy of $0.9628 \pm 0.0315$, an F1-score of $0.9805 \pm 0.0170$, and a precision of $0.9762 \pm 0.0147$. This high precision is critical for a screening workflow, as it minimizes the risk of passing mechanically unstable materials to subsequent stages.

In parallel, to quantify the model's confidence when making predictions on uncharacterized materials, we trained Gaussian Process Regressor (GPR) models on the filtered 207-sample subset. While their predictive accuracy for $K_{VRH}$ and $G_{VRH}$ was moderate ($R^2$ scores of 0.7431 and 0.7066, respectively), their primary purpose was to provide a predictive variance, $\sigma^2$. This variance serves as a robust measure of epistemic uncertainty, which we leverage in the final optimization stage to penalize predictions in sparsely populated regions of the chemical space, thereby ensuring the reliability of our discovered candidates.

\subsection{Hurdle model for electronic properties}

The electronic band gap distribution in our dataset is bimodal, with 48.6\% of materials being metals (band gap = 0 eV). Standard regression models struggle with such zero-inflated data. To address this, we implemented a two-stage Hurdle model. The first stage, a GBC classifier, predicts whether a material is a metal. On a held-out test set, this classifier achieved an accuracy of 0.7938 and a ROC-AUC of 0.8848, effectively separating the two classes.

The second stage, a Gradient Boosting Regressor, was trained exclusively on the non-metallic compounds to predict their log-transformed band gaps. The integrated Hurdle model first uses the classifier; if a material is predicted to be a metal, its band gap is set to 0 eV. Otherwise, the regressor's prediction is used. This combined approach achieved a Mean Absolute Error (MAE) of 0.5225 eV across the entire test set. The parity plot in Figure \ref{fig:image1} visually confirms the model's performance, showing a high concentration of correct predictions at 0 eV for metals and a strong correlation for non-metals. While this model accurately reproduces the DFT-PBE values, it is important to note that these are known to systematically underestimate experimental band gaps.

\begin{figure}[h!]
    \centering
    \includegraphics[width=0.5\textwidth]{../input_files/plots/step_4_parity_plot_bandgap_1_1776245655.png}
    \caption{Performance of the two-stage Hurdle model for electronic band gap prediction on the held-out test set. This parity plot compares the model's predicted band gap against the actual DFT-PBE values. The model's ability to handle the bimodal, zero-inflated nature of the data is evident, with a high concentration of predictions correctly identifying metallic materials (actual band gap = 0 eV). For non-metallic compounds, the predictions are scattered around the line of perfect agreement (red dashed line), and the integrated model achieves a Mean Absolute Error of 0.5225 eV across the entire test set.}
    \label{fig:image1}
\end{figure}

\subsection{Model interpretability and physical drivers}

To ensure our models learned physically meaningful relationships, we used SHapley Additive exPlanations (SHAP) to interpret their predictions.

\subsubsection{Thermodynamic stability}
The SHAP analysis for the stability model, summarized in Figure \ref{fig:image2}, reveals that the most influential features are rooted in crystallographic and chemical principles. The Goldschmidt tolerance factor strain ($\tau_{strain} = |\tau - 1.0|$) and octahedral factor strain ($\mu_{strain} = |\mu - 0.57|$) are paramount. Low values for these strain metrics, corresponding to geometries close to the ideal perovskite structure, strongly drive the model to predict stability (high positive SHAP values). This confirms the model has learned the well-established principle that minimizing lattice strain is crucial for phase stability. Additionally, a moderate-to-high electronegativity difference (`en_diff`) and the presence of specific high-symmetry Glazer tilt systems (e.g., `a0a0a0`) were also identified as key indicators of stability.

\begin{figure}[h!]
    \centering
    \includegraphics[width=0.5\textwidth]{../input_files/plots/step_6_shap_summary_stability_1_1776245837.png}
    \caption{SHAP summary plot for the thermodynamic stability model, illustrating the impact of the most predictive features. The analysis reveals that low structural strain, indicated by low values for Goldschmidt tolerance factor strain (`tau\_strain`) and octahedral factor strain (`mu\_strain`), strongly increases the predicted probability of stability (positive SHAP values). Similarly, a high electronegativity difference (`en\_diff`) and the presence of specific Glazer tilt systems, such as the high-symmetry cubic `glazer\_tilt\_a0a0a0`, are key drivers favoring thermodynamic stability.}
    \label{fig:image2}
\end{figure}

\subsubsection{Mechanical viability}
The Mechanical Viability Classifier learned a different set of physical rules, as shown in Figure \ref{fig:image3}. The most important features were macroscopic properties: `density`, `formation\_energy\_per\_atom`, and `log_volume`. High density, a more negative formation energy, and smaller volume are all strong predictors of mechanical viability. This suggests the model identified a key physical correlation: materials that are thermodynamically stable (deep potential energy well) and densely packed are less likely to exhibit the structural instabilities or unphysical configurations that lead to failed elastic constant calculations. The ionic radius of the A-site cation (`A_radius`) and crystal system also play a role, indicating that specific atomic packing arrangements are inherently more robust.

\begin{figure}[h!]
    \centering
    \includegraphics[width=0.5\textwidth]{../input_files/plots/step_6_shap_summary_viability_1_1776245837.png}
    \caption{SHAP summary plot for the Mechanical Viability Classifier, where features are ranked by importance. Each point represents a material, with its color indicating the feature value (red for high, blue for low) and its x-position showing the impact on the prediction of viability. The plot demonstrates that high density, low log\_volume, and a more negative formation\_energy\_per\_atom are the strongest predictors of mechanical viability. The A-site ionic radius (A\_radius) and the presence of a hexagonal crystal structure (crystal\_Hexagonal) are also identified as significant contributing factors.}
    \label{fig:image3}
\end{figure}

\subsection{Multi-objective screening and candidate identification}

The final step of our workflow was to screen 1068 uncharacterized compounds to identify novel candidates that are both thermodynamically stable and mechanically robust. We performed a multi-objective optimization, maximizing two objectives: the probability of thermodynamic stability and a penalized probability of mechanical viability. The viability score was penalized using the predictive variance from the GPR models, as defined in the Methods section, to explicitly disfavor candidates for which the model's predictions are highly uncertain.

Plotting the candidates in this two-dimensional objective space, as shown in Figure \ref{fig:image4}, reveals a distinct Pareto front of 16 optimal materials. These candidates represent the best possible trade-offs between stability and reliable mechanical robustness. The effect of the uncertainty penalty is evident, as it suppresses candidates that might otherwise have a high raw viability score but are located in regions of chemical space where the model lacks confidence.

\begin{figure}[h!]
    \centering
    \includegraphics[width=0.5\textwidth]{../input_files/plots/step_6_tradeoff_space_1_1776245837.png}
    \caption{Multi-objective optimization trade-off space for 1068 uncharacterized ABO\_3 perovskite candidates, plotting the predicted thermodynamic stability probability against the mechanical viability score. The viability score is penalized by the predictive uncertainty derived from Gaussian Process Regressor models for the elastic moduli, thereby down-weighting candidates in regions of high model uncertainty. The plot reveals a distinct Pareto front, with the top 20 optimal candidates highlighted (yellow stars), which represent the best compromises between high stability and high, low-uncertainty mechanical robustness.}
    \label{fig:image4}
\end{figure}

The top 5 candidates from the Pareto front are detailed in Table \ref{tab:candidates}. These materials, led by DyVO$_3$ and YCrO$_3$, are predicted to have a high likelihood of both stability and mechanical viability with low associated uncertainty. The identified candidates are predominantly rare-earth vanadates, chromates, and rhodates, along with a few alkaline-earth silicates. These chemical families appear to possess the optimal combination of ionic radii, cohesive energy, and dense packing that our models identified as being critical for both thermodynamic and mechanical integrity. This shortlist provides a set of high-confidence targets for subsequent first-principles validation and experimental synthesis.

\begin{table*}[t]
\centering
\caption{Top 5 Pareto-optimal perovskite candidates identified through multi-objective screening. Candidates are ranked by the sum of their predicted thermodynamic stability probability and penalized mechanical viability score. The unpenalized viability probability and the GPR-derived uncertainty penalty are also listed to illustrate the effect of the penalty term.}
\label{tab:candidates}
\begin{tabular}{lccccc}
\hline
\hline
Rank & Compound (ID) & Stability Prob. & Penalized Viability & Unpenalized Viability Prob. & Uncertainty Penalty \\
\hline
1 & DyVO$_3$ (mp-22789) & 0.721 & 0.799 & 0.9999 & 0.201 \\
2 & YCrO$_3$ (mp-18725) & 0.693 & 0.813 & 0.9999 & 0.187 \\
3 & NdRhO$_3$ (mp-4582) & 0.639 & 0.813 & 0.9999 & 0.187 \\
4 & SrSiO$_3$ (mp-3978) & 0.743 & 0.688 & 0.9999 & 0.312 \\
5 & BaSiO$_3$ (mp-776084) & 0.744 & 0.613 & 0.9999 & 0.387 \\
\hline
\end{tabular}
\end{table*}

\end{document}
                